% GENERATED FILE. Edit canonical lessons, then run npm run generate:books.
% canonical-source-hash: fnv1a64:0247278dffd94d47
% canonical-lessons: UR-W11-a2-connected-composition

\chapter{A connected A2 message}
\label{ch:ur-a-connected-a2-message}
\hlchaptermodality{143}

\begin{chapteropening}
\textbf{By the end of this chapter:} \emph{I can plan and write a connected 70–90 word Urdu message about a familiar day for a named reader and purpose.}

Complete chapter 143 by writing a model-free 70–90 word connected Urdu message for a named reader and practical purpose.
\end{chapteropening}

\section[(connected A2 writing)]{A connected message about your day}

\label{lesson:UR-W11-a2-connected-composition}



\begin{quote}
A connected message is not one long sentence. It is a short route for one reader. Before writing, draw eight empty boxes and give them these jobs:

\begin{enumerate}
\raggedright
  \item reader and greeting;
  \item reason for writing;
  \item where your day begins;
  \item two things you usually do;
  \item what happens next;
  \item one contrast or reason;
  \item one question for the reader;
  \item a polite close.
\end{enumerate}

Write only one or two English cue words in each box. Do not draft Urdu yet.
\end{quote}

\subsection*{Writing — connected composition}
There is no model, sentence bank, or copyable answer on this page.

\textbf{Reader:} Sana. \textbf{Purpose:} help Sana understand your usual weekday and choose a good time to visit you tomorrow.

Write \textbf{70–90 words in Urdu script}. Use every planning box. Include at least four familiar actions, two time anchors, one preference, one contrast or reason, one sequence link, and one direct question for Sana. Keep the reader visible by making every detail help with the visit. The instructions give content, not sentences: choose and connect the Urdu yourself.

This first connected attempt is untimed. Put away Roman Urdu, Devanagari, dictionaries, translators, and spell-checkers. Either approved Nastaliq or accessibility Naskh is fine. Ordinary legible Urdu handwriting is enough; imitating Nastaliq calligraphic slope or ligatures is not part of the task.

When the whole message exists, count its words once. If it is short, add one useful detail to a planning box. If it is long, remove a detail that does not help Sana. Do not replace the message with a memorised model.

\begin{tcolorbox}[breakable,colback=teal!4,colframe=teal!35!black,title={Before you move on}]
Review the same attempt in four separate passes:

\begin{enumerate}
\raggedright
  \item \textbf{Task fulfilment and relevant content} — Sana, the weekday, tomorrow's visit, and every requested content point are present.
  \item \textbf{Comprehensibility and organisation} — the eight jobs form a route rather than a list of disconnected sentences.
  \item \textbf{Urdu vocabulary, grammar, and register} — connectors make addition, sequence, contrast, and reason clear in a polite everyday tone.
  \item \textbf{Urdu orthographic control} — letters, dots, joins, non-joins, word boundaries, and punctuation remain readable. Unwritten short vowels and optional diacritics receive neither bonus nor penalty unless they change the intended word.
\end{enumerate}

Mark one strength and one next step. Repair only the next step; keep the original attempt so the next 70–90 word message can show real growth.
\end{tcolorbox}
